\chapter{Introduction}
\label{ch:intro}

Circe is a verification pipeline from C, and from STL-free C++-lite, to
Lean~4. Source programs travel through ClangIR (CIR, the raw
\texttt{CIRGen} output) into pure, memory-free Lean~4 definitions produced
by a verified emitter, and functional-correctness proofs are stated as pure
equations over those definitions. The style is Aeneas-like: instead of
reasoning about a heap with separation logic, the pipeline translates
programs into total-looking value functions and proves things about the
functions~\cite{aeneas-site}.

Two design bets carry the whole project. First, the common case needs no
memory model at all: admitted programs are value programs, so their proofs
are equations, inductions, and tactic automation rather than framing
lemmas. Second, whatever falls outside the admitted subset is rejected
loudly with a dedicated message. Circe never silently models memory it
does not understand; an out-of-subset program is an error, not an
approximation.

This report is the long-term home for the ideas behind the pipeline: how
C/C++ is interpreted and given semantics in Lean
(Chapter~\ref{ch:semantics}), what the verification rests on
(Chapter~\ref{ch:verification}), how the repository and harness fit
together (Chapter~\ref{ch:technical}), where the work sits among related
efforts (Chapter~\ref{ch:related}), and where the codebase stands
(Chapter~\ref{ch:progress}). The Markdown documents under
\texttt{docs/} and the Lean sources themselves remain normative; this
report explains them and points at them rather than duplicating them.

A note on freshness: only Chapter~\ref{ch:progress} tracks the moving
codebase, and it does so through regenerated snapshots, not hand edits.
Every other chapter describes stable designs and changes only when the
design itself changes. The snapshot built into this copy is identified by
the git SHA printed in Chapter~\ref{ch:progress}.
